\documentclass[11pt]{article}
\usepackage{amsmath, amsthm, amssymb}
\usepackage{array}
\usepackage{forest}
\DeclareMathOperator*{\argmax}{arg\,max}
\DeclareMathOperator*{\argmin}{arg\,min}
\usepackage[margin=1in]{geometry}
\theoremstyle{definition}
\newtheorem{definition}{Definition}[section]
\theoremstyle{plain}
\newtheorem{theorem}{Theorem}[section]
\title{}
\date{}
\begin{document}
\section{Definitions}
As usual, let $[N]=\{1,...,N\}$ be a finite set of agents. Imagine a single type of
good is being auctioned off and each agent $i \in [N]$ has a private value $v_i \in
\mathbb {R}_{\geq 0}$ they assign per unit of the good.
\\
Let $\mathbb{B}=\mathbb{R}_{\geq 0}^N$ be the set of possible bid profiles. Given
$b \in \mathbb{B}$ and $i \in [N]$ write $b_i$ to denote the $i$th component of $b$
and $b_{-i}$ to denote the remaining components. Note $v=(v_i)_{i \in [N]}\in \mathbb{B}$.
\\
Let $\mathbb{A}=\mathbb{R}_{\geq 0}^N$ be the set of possible allocations for how
much of the good each agent receives.
\\
Let $\mathbb{P}=\mathbb{R}_{\geq 0}^N$ be the set of possible payments for how much
each agent pays the auctioneer.
\\
\\
An allocation rule is a function $x: \mathbb{B} \to \mathbb{A}$. The goal of an
auctioneer is to find out $v \in \mathbb{B}$ and allocate resources according to
$x(v)$. We write $x(b_i;b'_{-i})$ to denote the output of $x$ when given an input
bid profile whose $i$th component is $b_i$ and whose other components are given by
$b'_{-i}.$ Additionally we write $x_i(b)$ to denote the $i$th component of $x(b)$,
i.e. how much of the good agent $i$ is allocated when everyone bids according to $b
\in \mathbb{B}.$ The set $x(\mathbb{B})=\{a \in \mathbb{A} : \exists b \in
\mathbb{B}, x(b)=a\}$ is often called the set of feasible allocations. An
allocation rule $x: \mathbb{B} \to \mathbb{A}$ is monotone iff $\forall i \in [N]$
and $\forall b,b' \in \mathbb{B}$ if $b'_i > b_i$ and $b'_{-i}=b_{-i}$ then
$x_i(b') \geq x_i(b)$.
\\
\\
A payment rule is a function $p: \mathbb{B} \to \mathbb{P}$. Again, we write
$p(b_i;b'_{-i})$ to denote the output of $p$ when given an input bid profile whose
$i$th component is $b_i$ and whose other components are given by $b'_{-i}$.
Additionally, we write $p_i(b)$ to denote the $i$th component of $p(b)$, i.e. how
much agent $i$ pays the auctioneer whenever everyone bids according to $b\in
\mathbb{B}.$ For all $i \in [N]$ we take agent $i$'s payoff when everyone bids $b
\in \mathbb{B}$ to be given by $v_i \cdot x_i(b)-p_i(b)$. The point of the payment
rule is to help the auctioneer find out $v \in \mathbb{B}$ so they may allocate
resources according to $x(v)$.
\\
\\
We say that a payment rule $p$ implements the allocation rule $x$ iff:
\\
1. (Strategy Proofness)- $\forall v \in \mathbb{B}, \forall i \in [N], \forall b
\in \mathbb{B}, v_i \cdot x_i(v_i;b_{-i})-p_i(v_i;b_{-i}) \geq v_i \cdot x_i(b)-
p_i(b)$.
\\
2. (Individual Rationality)- $\forall v \in \mathbb{B}, \forall i \in [N], \forall
b \in \mathbb{B}, v_i \cdot x_i(v_i;b_{-i})-p_i(v_i;b_{-i}) \geq 0$.
\\
Strategy proofness ensures an always best response for each agent is to just bid
their true valuation. Individual rationality ensures that bidding their true
valuation always results in nonnegative utility. Together these conditions imply
that, by using the payment rule $p$ along with the allocation rule $x$, everyone is
incentivized to participate in the auction truthfully and the final allocation will
be given by the allocation $x(v) \in \mathbb{A}.$
\section{Examples}
\subsection*{Single Indivisible Unit to Highest Bidder}
Consider the allocation rule $x:\mathbb{B} \to \mathbb{A}$ which $\forall b \in
\mathbb{B}$ and $\forall i \in [N]$ sets:
$$x_i(b)=\begin{cases}1 & i=\min\argmax_{j \in [N]} b_j \\ 0 & \text{else}.
\end{cases}$$
Note $x(\mathbb{B})$ consists of all tuples in $\mathbb{R}^N_{\geq 0}$ with exactly
1 entry equal to 1 and the remaining $N-1$ entries equal to $0$.
\subsection*{$k$ Indivisible Units to $k$ Highest Bidders}
Suppose $k< N$. Consider the allocation rule $x:\mathbb{B} \to \mathbb{A}$ which $
\forall b \in \mathbb{B}$ and $\forall i \in [N]$ sets:
$$x_i(b)=\begin{cases}1 & i \in \lambda(b) \\ 0 & \text{else}. \end{cases}$$
where $\lambda(b)$ is the set of $k$ agents such that $\forall i \in \lambda(b),
\forall j \in [N] \backslash \lambda(b)$ we have $b_i \geq b_j$ and if $b_i=b_j$ we
have $i < j$. Note $x(\mathbb{B})$ consists of all tuples in $\mathbb{R}^N_{\geq
0}$ with exactly $k$ entries equal to 1 and the remaining $N-k$ entries equal to
$0$.
\subsection*{Single Divisible Unit Allocated Proportional to Bid}
Consider the allocation rule $x:\mathbb{B} \to \mathbb{A}$ which $\forall b \in
\mathbb{B}$ and $\forall i \in [N]$ sets:
$$x_i(b)=\begin{cases} \frac{b_i}{\sum_{j \in [N]} b_j} & b_i \neq 0 \\ 0 &
\text{else}. \end{cases}$$
Note $x(\mathbb{B})$ consists of all tuples in $\mathbb{R}^N_{\geq 0}$ whose
entries sum to $1$ as well as the tuple of all $0$'s.
\subsection*{First Price vs. Second Price}
Consider the allocation rule $x: \mathbb{B} \to \mathbb{A}$ which allocates a
single indivisible unit to the highest bidder (our first example). Consider the
payment rules $p^1: \mathbb{B} \to \mathbb{P}$ and $p^2: \mathbb{B} \to \mathbb{P}$
which $\forall b \in \mathbb{B}$ and $\forall i \in [N]$ set:
$$p^1_i(b)=\begin{cases}b_i & i=\min \argmax_{j \in [N]} b_j \\ 0 & \text{else}
\end{cases}$$
$$p^2_i(b)=\begin{cases} \max_{j \in [N] \backslash \{i\}} b_j & i=\min \argmax_{j
\in [N]} b_j \\ 0 & \text{else} \end{cases}.$$
Under $p^1$, the agent who receives the item under $x$ pays the amount they bid.
While $p^1$ satisfies individual rationality for $x$ it does not satisfy strategy
proofness. Under $p^2$, the agent who receives the item under $x$ pays the highest
bid amongst everyone who does not receive the item under $x$. Note $p^2$
successfully implements $x$ as it satisfies both individual rationality and
strategy proofness.
\section{Theorems}
\noindent\textbf{Theorem 19:} There exists a payment rule $p: \mathbb{B} \to
\mathbb{P}$ which implements the allocation rule $x: \mathbb{B} \to \mathbb{A}$ iff
$x$ is monotone. In particular, for any monotone allocation rule $x: \mathbb{B} \to
\mathbb{A}$ there exists a unique payment rule $p: \mathbb{B} \to \mathbb{P}$ which
implements it.
\\
\\
\underline{Proof}
\\
Suppose there exists a payment rule $p: \mathbb{B} \to \mathbb{P}$ which implements
the allocation rule $x: \mathbb{B} \to \mathbb{A}$. Take arbitrary $i \in [N]$ and
$b,b' \in \mathbb{B}$ such that $b'_i > b_i$ and $b'_{-i}=b_{-i}$. Imagine agents'
true valuations are given by $v=b' \in \mathbb{B}$. Then strategy proofness
implies:
$$b'_i \cdot x_i(b')-p_i(b') =b'_i \cdot x_i(b'_i;b_{-i})-p_i(b'_i;b_{-i}) \geq
b'_i \cdot x_i(b)-p_i(b).$$
Rearranging, this means $b'_i \cdot (x_i(b')-x_i(b)) \geq p_i(b')-p_i(b)$.
\newpage
\noindent Now imagine agent's true valuations are given by $v=b \in \mathbb{B}$.
Similar to before, strategy proofness implies:
$$b_i \cdot x_i(b) - p_i(b)=b_i \cdot x_i(b_i; b'_{-i})-p_i(b_i; b'_{-i}) \geq b_i
\cdot x_i(b')-p_i(b').$$
Rearranging, this means $p_i(b')-p_i(b) \geq b_i \cdot (x_i(b')-x_i(b))$. Hence:
$$b'_i \cdot (x_i(b')-x_i(b)) \geq p_i(b')-p_i(b) \geq b_i \cdot (x_i(b')-x_i(b)) $$
and so:
$$(b'_i-b_i) \cdot (x_i(b')-x_i(b)) \geq 0.$$
Since $b'_i -b_i>0$, we must have $x_i(b')-x_i(b) \geq 0$, i.e. $x_i(b') \geq
x_i(b)$. Thus $x$ is monotone.
\\
\\
Now suppose $x: \mathbb{B} \to \mathbb{A}$ is monotone. Take arbitrary $i \in [N]$
and $y,z \in \mathbb{R}_{\geq 0}$ such that $y < z$. Additionally, take arbitrary
$b \in \mathbb{B}$. If $x$ were to be implementable by some payment rule $p:
\mathbb{B} \to \mathbb{P}$, the preceding argument shows us:
$$y \cdot (x_i(z;b_{-i})-x_i(y;b_{-i})) \leq p_i(z;b_{-i})-p_i(y;b_{-i}) \leq z
\cdot (x_i(z;b_{-i})-x_i(y;b_{-i})).$$
In particular, since $y$ and $z$ are arbitrary, we can take $(r_t)_{t \in
\{0,...,m\}}$ to be given by $r_t=b_i \cdot \frac{t}{m}$ with $m>1$ so $r_0=0$
while $r_m=b_i$ and conclude that:
$$\sum_{t \in [m]} r_{t-1} \cdot (x_i(r_t;b_{-i})-x_i(r_{t-1};b_{-i})) \leq
p_i(b)-p_i(0;b_{-i}) \leq \sum_{t \in [m]} r_t \cdot (x_i(r_t;b_{-i})-
x_i(r_{t-1};b_{-i})).$$
since $\sum_{t \in [m]} (p_i(r_t;b_{-i})-p_i(r_{t-1};b_{-i}))=p_i(r_m;b_{-i})-
p_i(r_0;b_{-i})=p_i(b)-p_i(0;b_{-i})$.
\\
\\
Rearranging we have:
\begin{align*}
\sum_{t \in [m]} r_{t-1} \cdot (x_i(r_t;b_{-i})-x_i(r_{t-1};b_{-i}))&=[r_{m-1}
\cdot x_i(r_m;b_{-i})-r_{0} \cdot x_i(r_0;b_{-i})]-\sum_{t \in [m-1]} x_i(r_t;b_{-
i}) \cdot (r_t-r_{t-1})
\\
&=r_{m-1} \cdot x_i(b)-\sum_{t \in [m-1]} x_i(r_t;b_{-i}) \cdot (r_t-r_{t-1})
\end{align*}
and:
\begin{align*}
\sum_{t \in [m]} r_t \cdot (x_i(r_t;b_{-i})-x_i(r_{t-1};b_{-i}))&=[r_{m} \cdot
x_i(r_m;b_{-i})-r_{1} \cdot x_i(r_0;b_{-i})]-\sum_{t \in [m-1]} x_i(r_t;b_{-i})
\cdot (r_{t+1}-r_{t})
\\
&=[b_i \cdot x_i(b)-r_1 \cdot x_i(0;b_{-i})]-\sum_{t \in [m-1]} x_i(r_t;b_{-i})
\cdot (r_{t+1}-r_{t})
\end{align*}
Since $x$ is monotone we know the function $x_i(\ \cdot \ ; b_{-i})$ is
nondecreasing and defined everywhere on the interval $[0,b_i]$ it follows $x_i(\
\cdot \ ; b_{-i})$ is Riemann integrable over the interval $[0,b_i]$. Hence:
$$\lim_{m \to \infty } \left(r_{m-1} \cdot x_i(b)-\sum_{t \in [m-1]} x_i(r_t;b_{-
i}) \cdot (r_t-r_{t-1})\right)=b_i \cdot x_i(b)-\int_{0}^{b_i} x_i(w;b_{-i}) \ dw$$
and:
$$\lim_{m \to \infty } \left([b_i \cdot x_i(b)-r_1 \cdot x_i(0;b_{-i})]-\sum_{t \in
[m-1]} x_i(r_t;b_{-i}) \cdot (r_{t+1}-r_{t})\right)=b_i \cdot x_i(b)-\int_{0}^{b_i}
x_i(w;b_{-i}) \ dw.$$
Since the first limit is a lower bound for $p_i(b)-p_i(0;b_{-i})$ and the second is
an upper bound it follows that:
$$p_i(b)-p_i(0;b_{-i})=b_i \cdot x_i(b)-\int_{0}^{b_i} x_i(w;b_{-i}) \ dw.$$
Finally, imagine $v \in \mathbb{B}$ is the bid profile where $v_i=0$ and $v_{-i}
=b_{-i}.$ Since we have assumed $p$ satisfies individual rationality for $x$, we
know $0 \cdot x_i(0;b_{-i})-p_i(0;b_{-i}) =-p_i(0;b_{-i})\geq 0$, i.e. $p_i(0;b_{-
i})=0$ as $p$ only outputs elements of $\mathbb{R}^N_{\geq 0}$.
\\
\\
Therefore, if $x$ were to be implementable, it must be implemented by the payment
rule $p: \mathbb{B} \to \mathbb{P}$ which $\forall b \in \mathbb{B}$ and $\forall i
\in [N]$ sets:
$$p_i(b)=b_i \cdot x_i(b)-\int_{0}^{b_i} x_i(w;b_{-i}) \ dw$$
and no other. We will now show that this payment rule indeed implements $x$.
\\
\\
Take arbitrary $v \in \mathbb{B}$, $i \in \mathbb{N}$, and $b \in \mathbb{B}$.
Note:
\begin{align*}
v_i\cdot x_i(v_i;b_{-i})-p_i(v_i;b_{-i})&=v_i\cdot x_i(v_i;b_{-i})-v_i \cdot
x_i(v_i;b_{-i})+\int_{0}^{v_i} x_i(w;b_{-i}) \ dw
\\
&=\int_{0}^{v_i} x_i(w;b_{-i}) \ dw
\\
&\geq 0
\end{align*}
since $x$ only outputs elements of $\mathbb{R}^N_{\geq 0}$. Thus $p$ satisfies
individual rationality for $x$. Moreover, note:
$$v_i\cdot x_i(v_i;b_{-i})-p_i(v_i;b_{-i})=\int_{0}^{v_i} x_i(w;b_{-i}) \ dw
=\int_{b_i}^{v_i}x_i(w;b_{-i}) \ dw+\int_{0}^{b_i}x_i(w;b_{-i}) \ dw$$
If $b_i \geq v_i$ then $\forall w \in [v_i,b_i]$ we know $x_i(b) \geq x_i(w;b_{-i})
$ by monotonicity and so:
$$\int_{v_i}^{b_i} x_i(b) \ dw \geq \int_{v_i}^{b_i} x_i(w; b_{-i}) \ dw.$$
If $b_i \leq v_i$ then $\forall w \in [b_i,v_i]$ we know $x_i(b) \leq x_i(w;b_{-i})
$ by monotonicity and so:
$$\int_{b_i}^{v_i} x_i(b) \ dw \leq \int_{b_i}^{v_i} x_i(w; b_{-i}) \ dw.$$
Hence in either case:
$$\int_{b_i}^{v_i} x_i(w; b_{-i}) \ dw \geq \int_{b_i}^{v_i} x_i(b) \ dw=v_i\cdot
x_i(b)-b_i\cdot x_i(b).$$
Thus we have:
\begin{align*}
v_i\cdot x_i(v_i;b_{-i})-p_i(v_i;b_{-i})&=\int_{b_i}^{v_i}x_i(w;b_{-i}) \ dw+
\int_{0}^{b_i}x_i(w;b_{-i}) \ dw
\\
&\geq v_i\cdot x_i(b)-b_i\cdot x_i(b)+\int_{0}^{b_i}x_i(w;b_{-i}) \ dw
\\
&=v_i\cdot x_i(b)-\left(b_i\cdot x_i(b)-\int_{0}^{b_i}x_i(w;b_{-i}) \ dw
\right)
\\
&=v_i\cdot x_i(b)-p_i(b)
\end{align*}
, i.e. $p$ satisfies strategy proofness for $x$ and therefore $p$ implements $x$.
$\square$
\end{document}
